% LuaLaTeX-ja による日本語Beamerスライド
\documentclass[unicode]{beamer}
\usepackage{luatexja}
\usepackage{luatexja-fontspec}
\usetheme{Madrid}

\title{日本語組版とBeamerスライド}
\subtitle{LuaTeX-ja環境におけるマルチバイト検証}
\author{開発チーム}
\institute{次世代エディタ研究所}
\date{2026年10月}

\begin{document}

\begin{frame}
  \titlepage
\end{frame}

\section{概要と背景}

\begin{frame}{日本語組版の課題と対策}
  \framesubtitle{全角文字とインライン記号の共存}
  日本語組版では、約物(「」『』【】・※★)や全角空白の混在に注意が必要です。

  \begin{itemize}
    \item 具象構文木(CST)による完全ロスレス表現
    \item マルチバイトUTF-8境界でのオフセット計算の整合性
    \item 日本語コメント行の改行コード保持(LF / CRLF)
  \end{itemize}

  \begin{block}{確認事項}
    パースおよび再文字列化において、1バイトのズレも生じないことを確認します。
  \end{block}
\end{frame}

\end{document}
